\section{Exact rigidity for chi-square quadratic forms}\label{sec:exact}

We prove Theorem~\ref{thm:exact-psd} and then apply it to the ordinary sample variance. Write $A=B^{\mathsf T}B$, where $B\in\R^{r\times d}$ has full row rank, and let $b_i$ be the $i$th column of $B$. Thus $A_{ii}=\|b_i\|^2$.

We first note that the chi-square hypothesis automatically gives the MGF regularity needed below. Fix an active coordinate $i$, so $A_{ii}>0$, and let
\[
Q:=X^{\mathsf T}AX\sim\chi_r^2.
\]
Independence and centering imply
\[
\E(Q\mid X_i)=A_{ii}X_i^2+\sum_{j\ne i}A_{jj}\sigma_j^2.
\]
For every $0<c<1/2$, conditional Jensen therefore yields
\[
\exp\!\left(cA_{ii}X_i^2+c\sum_{j\ne i}A_{jj}\sigma_j^2\right)
\le \E(e^{cQ}\mid X_i).
\]
Taking expectations and using the chi-square MGF gives
\begin{equation}\label{eq:auto-square-exp}
\E e^{cA_{ii}X_i^2}<\infty.
\end{equation}
Consequently $M_i(s)<\infty$ for every $s\in\R$, by Young's inequality
\[
sx\le cA_{ii}x^2+\frac{s^2}{4cA_{ii}}.
\]

If $V$ is uniform on $S^{r-1}$, write
\[
\Psi_r(z):=\E_V e^{zV_1}.
\]
The transform is well defined for all $z\in\C$ because $V_1$ is bounded. Spherical averaging gives
\begin{equation}\label{eq:spherical-transform}
\E_V\prod_{i=1}^d M_i\bigl(t\langle b_i,V\rangle\bigr)
 =\E_{X,V}e^{t\langle BX,V\rangle}
 =\E\Psi_r\bigl(t\sqrt{X^{\mathsf T}AX}\bigr).
\end{equation}
Since $X^{\mathsf T}AX\sim\chi_r^2$, the final quantity in \eqref{eq:spherical-transform} equals
\begin{equation}\label{eq:spherical-gaussian}
\E_{G,V}e^{t\langle G,V\rangle}=e^{t^2/2},
\qquad G\sim\Normal(0,I_r).
\end{equation}
Only active coordinates enter the product, because $A_{ii}=0$ implies $b_i=0$.

For each active coordinate, let the local dominance hold on $|s|<\rho_i$. Since there are finitely many active coordinates, reducing to a common radius $\rho_*:=\min_i\rho_i>0$ makes all subsequent expressions valid for sufficiently small $|t|$.

Define
\[
C_i(s):=\log M_i(s)-\frac{\sigma_i^2s^2}{2}
\]
and
\begin{equation}\label{eq:Ht-exact}
H_t(v):=\sum_{i=1}^d C_i\bigl(t\langle b_i,v\rangle\bigr)
 +\frac{t^2}{2}\left(\sum_{i=1}^d\sigma_i^2\langle b_i,v\rangle^2-1\right).
\end{equation}
The quadratic terms cancel inside the exponential, so \eqref{eq:spherical-gaussian} implies
\[
\E_V e^{H_t(V)}=e^{-t^2/2}\E_V\prod_iM_i(t\langle b_i,V\rangle)=1.
\]
Jensen's inequality therefore gives $\E_VH_t(V)\le0$.

The mean of the quadratic correction in \eqref{eq:Ht-exact} is zero. Indeed,
\[
\E_V\sum_i\sigma_i^2\langle b_i,V\rangle^2
 =\frac1r\sum_i\sigma_i^2\|b_i\|^2
 =\frac1r\E(X^{\mathsf T}AX)=1,
\]
where the last equality uses the mean of $\chi_r^2$. Since $V$ and $-V$ have the same law, coordinatewise dominance now gives
\begin{align*}
2\E_VH_t(V)
 &=\sum_{i=1}^d\E_V\left[
 \log M_i(t\langle b_i,V\rangle)+\log M_i(-t\langle b_i,V\rangle)
 -\sigma_i^2t^2\langle b_i,V\rangle^2\right]\\
 &=\sum_{i=1}^d\E_V R_i\bigl(t\langle b_i,V\rangle\bigr)\ge0,
\end{align*}
where $R_i(s)=\log M_i(s)+\log M_i(-s)-\sigma_i^2s^2$. Hence $\E_VH_t(V)=0$, and every nonnegative summand in the last display has zero expectation.

For an active coordinate, $A_{ii}=\|b_i\|^2>0$. If $r\ge2$, the random variable $\langle b_i,V\rangle$ has a density that is strictly positive on the interior of $[-\|b_i\|,\|b_i\|]$. Continuity of $R_i$ therefore implies $R_i(s)=0$ on a neighborhood of the origin. If $r=1$, then $V\in\{-1,1\}$ and varying $t$ directly gives the same conclusion, because $R_i$ is even. Thus, for every active coordinate,
\[
M_i(s)M_i(-s)=e^{\sigma_i^2s^2}
\]
near the origin. If $\sigma_i^2=0$, then $X_i=0$ almost surely. If $\sigma_i^2>0$, an independent copy $X_i'$ gives that $X_i-X_i'$ has the MGF of $\Normal(0,2\sigma_i^2)$ near zero and hence has that Gaussian law. The L\'evy--Cram\'er decomposition theorem applied to the independent sum $X_i+(-X_i')$ then implies that $X_i$ is Gaussian.

This proves Theorem~\ref{thm:exact-psd}. For Corollary~\ref{cor:exact-sample-var}, take
\[
A=I_n-\frac1n\mathbf 1\mathbf 1^{\mathsf T}.
\]
It is a rank-$n-1$ orthogonal projection, every diagonal entry is $(n-1)/n>0$, and its quadratic form is exactly $\sum_i(X_i-\bar X)^2$. The corollary follows immediately.
